\documentclass[11pt]{article}
\usepackage[margin=1in]{geometry}
\usepackage[T1]{fontenc}
\usepackage{lmodern}
\usepackage{xcolor}
\usepackage{parskip}
\definecolor{ink}{HTML}{17231F}
\definecolor{orange}{HTML}{D9653B}
\begin{document}
\color{ink}
{\Large\bfseries Generative AI Portfolio Reflection\\[4pt]}
{\color{orange}\small Truong Huu Loi \textbar\ 張友利 \textbar\ M1561023 \textbar\ HW\#1}

\vspace{12pt}
I used generative AI as a collaborator for both the planning and implementation of my portfolio website. I first asked AI to interpret the assignment and convert it into practical requirements: four clear sections, responsive behavior, a meaningful visual element, a public-facing contact area, an AI interaction log, and a reflection. This helped me focus on the experience of a professor, employer, or collaborator who should understand my background within about two minutes.

I then provided my CV and asked AI to turn it into concise portfolio copy. The result accurately highlighted my work as a Process Data Analyst at JAPFA Comfeed Vietnam, my Data Engineer internship at ESTEC Vietnam, and my Data Science education at Ho Chi Minh City University of Industry. It also helped explain my research project about FOMC decisions, VND interest rates, and exchange rates using machine learning and natural language processing. I added that I am currently studying at CGU, because this was current information that was not included in the older CV.

AI was especially useful for generating a visual direction and a first implementation quickly. It proposed an editorial layout with a calm cream background, dark text, a lime accent, oversized headings, a timeline, and an analytical trend graphic. I chose a CSS/SVG graph rather than a stock image so the visual would be lightweight, original, and connected to my field. I also asked AI to review the page for mobile behavior, link functionality, accessibility, privacy, and assignment compliance.

I did not accept the output blindly. I shortened several paragraphs so the site would be easier to scan, kept only claims supported by my CV, and used “CGU” without guessing its full institutional name. I also kept the Chinese-name spelling easy to revise because it should be confirmed rather than invented. I checked that the navigation anchors, email link, and phone link matched the visible content and that the layout changes to a single-column structure on small screens.

This process taught me that AI is strong at generating alternatives, structure, and implementation speed, but human judgment is necessary for accuracy, privacy, tone, and prioritization. The best result came from treating AI output as a draft: I supplied constraints, inspected the claims and interaction details, and made decisions based on my actual experience and the assignment audience.
\end{document}
